The Schwartzberg score and Polsby-Popper score are algebraically equivalent:
$S = 1/\sqrt{\text{PP}}$ exactly (Theorem~\ref{thm:equivalence}).
Despite this equivalence, the two representations serve different legislative
and litigation contexts.
PP's $[0,1]$ scale is standard in academic literature and federal court filings;
S's $[1,\infty)$ scale is preferred by Colorado's IRC and several state
redistricting commissions.

Our empirical results confirm that:
\begin{enumerate}
  \item Ratio-optimal achieves mean S $\approx 2.13^\dagger$ on NC 2020
    (mean PP $\approx 0.22^\dagger$), the best performance among all four algorithms.
  \item Colorado ratio-optimal results achieve mean S $\approx 1.74^\dagger$
    with maximum district S $= 2.00$, satisfying any $S \leq 2.0$ threshold
    across all 8 districts.
  \item The conversion table (Section~4) enables practitioners to translate
    any bisect PP output to S without additional computation.
\end{enumerate}

The practical recommendation is to report PP as the primary metric with S
as a parenthetical conversion in jurisdictions that mandate Schwartzberg-equivalent
criteria.
The K.7 composite guide provides template language incorporating both representations.
